\section{Alcance}
El presente procedimiento establece los lineamientos formales que SKYFILMS aplicará para la evaluación, selección, aprobación y control continuo de los proveedores u organizaciones externas que presten servicios de mantenimiento (menor, mayor y correctivo) a la flota de Sistemas de Aeronaves No Tripuladas (UAS) de la compañía. El objetivo es garantizar que cualquier intervención técnica externalizada cumpla con los estándares de aeronavegabilidad, el fabricante y los lineamientos de Seguridad Operacional (SMS) del RAC 100.

\section{Criterios y Requisitos para la Selección de Proveedores}
Para que un proveedor externo sea considerado y aceptado por SKYFILMS para realizar labores de mantenimiento a nuestras aeronaves, deberá cumplir satisfactoriamente con los siguientes criterios de evaluación:

\begin{enumerate}
    \item \textbf{Capacidad Técnica y de Infraestructura:} Demostrar que cuenta con instalaciones adecuadas, control ambiental (si aplica para componentes electrónicos/baterías), y áreas segregadas para la recepción, intervención y almacenamiento de los UAS y sus componentes.

    \item \textbf{Herramientas e Insumos:} Poseer herramientas calibradas y adecuadas (con trazabilidad metrológica), así como acceso a insumos y repuestos originales o aprobados por el fabricante del UAS.

    \item \textbf{Competencia del Personal:} Contar con personal técnico debidamente capacitado, con experiencia demostrable en la marca y modelo específico de los UAS operados por SKYFILMS, y conocer los manuales técnicos vigentes del fabricante.

    \item \textbf{Documentación y Trazabilidad:} Evidenciar un sistema de control documental que garantice el correcto registro del trabajo realizado (diligenciamiento de bitácoras, registros de mantenimiento, tarjetas de trabajo y emisión de certificados de conformidad o retorno al servicio).

    \item \textbf{Cultura de Seguridad (SMS):} Demostrar compromiso con prácticas de reporte de peligros y gestión de riesgos, alineadas con los principios de Seguridad Operacional que promueve SKYFILMS
\end{enumerate}

\section{Procedimiento de Selección y Control Continuo}
El Gerente de Seguridad Operacional y/o el Tecnico UAS de SKYFILMS ejecutarán el siguiente ciclo para garantizar la calidad del servicio tercerizado:

\begin{enumerate}
    \item \textbf{Identificación y Solicitud:} Ante la necesidad de tercerizar un mantenimiento, se solicitará al posible proveedor su portafolio de servicios, manuales de procedimientos técnicos y certificados de su personal.

    \item \textbf{Auditoría de Evaluación Inicial:} Se realizará una verificación documental y, de ser necesario, una visita a las instalaciones del proveedor para auditar su infraestructura, herramientas, control de inventarios y banco de pruebas.

    \item \textbf{Aprobación del Proveedor:} Si la evaluación es satisfactoria y cumple con los criterios de la sección \textit{Criterios y Requisitos para la Selección de Proveedores} de este capítulo, el proveedor será ingresado a la "Lista de Organizaciones Aceptadas" de SKYFILMS.

    \item \textbf{Control Continuo y Evaluación de Desempeño:} SKYFILMS monitoreará la calidad del mantenimiento mediante la revisión estricta de cada UAS tras su retorno al servicio. Adicionalmente, se realizarán auditorías recurrentes (anuales o tras la ocurrencia de un evento de seguridad operacional) para verificar que el proveedor mantiene los estándares por los cuales fue seleccionado. Si se detectan no conformidades, el proveedor podrá ser suspendido de la lista de organizaciones aceptadas hasta que se implementen acciones correctivas.
\end{enumerate}

\textbf{NOTA:} La trazabilidad del control de los proveedores se registra en el \textbf{Anexo 2 - Seguimiento Mantenimiento Taller}.

\section{Contrato, Representante Técnico y Auditoría Periódica}
\begin{itemize}
    \item Con cada proveedor aceptado se firma un contrato o acuerdo escrito que dice qué trabajos hace, que usa los manuales del fabricante, que entrega informe y certificado por cada trabajo, y que permite a Skyfilms y a la Aerocivil inspeccionar sus instalaciones. Se guarda una copia en el archivo técnico.
    \item El proveedor designa un representante técnico que firma los informes y responde por el trabajo.
    \item Antes de aceptarlo y luego cada año, el Gerente SMS y el Técnico UAS auditan al proveedor con una lista de chequeo (personal, herramientas, instalaciones, registros) y anotan el resultado en el Anexo 2 (F-018). Si la auditoría falla, se retira de la lista hasta corregir.
\end{itemize}

\section{Organizaciones y Lugares Aceptados para el Mantenimiento}
Tras la aplicación rigurosa de nuestro procedimiento de selección, SKYFILMS ha evaluado, inspeccionado y aceptado a las siguientes organizaciones e instalaciones para llevar a cabo el mantenimiento de nuestra flota UAS:

\begin{itemize}
    \item \textbf{Nombre de la Organización Proveedora:} [EN PROCESO DE SELECCIÓN Y APROBACIÓN]

    \item \textbf{Ubicación de las Instalaciones (Lugar de Mantenimiento):} [EN PROCESO DE SELECCIÓN Y APROBACIÓN]

    \item \textbf{Alcance del Mantenimiento Autorizado:} [EN PROCESO DE SELECCIÓN Y APROBACIÓN]
\end{itemize}

\textbf{Mientras no exista una organización aceptada:} Skyfilms no externaliza mantenimiento. Las fallas que superan los límites del Técnico UAS dejan la aeronave NO AERONAVEGABLE hasta que se acepte un proveedor siguiendo el procedimiento de este capítulo, que puede ser un centro de servicio autorizado del fabricante. Cuando se acepte, se actualizan esta sección, el Manual de Operaciones y la Aerocivil se informa como cambio sustancial.
